Take
for the family @@M0@@ the family of all elements of @@M1@@. The hypotheses of 1.2 are then satisfied, and consequently the homomorphism:
@@M2@@
is surjective. By the definition of @@M3@@, this is equivalent to saying that the homomorphism:
@@M4@@
is surjective.

\sgaStableTarget{sga6.e09.cor.1-4}\textbf{Corollary 1.4.} \uline{Let} @@M5@@ \uline{be a flat morphism of noetherian schemes}. \uline{Suppose that for every point} @@M6@@ \uline{of} @@M7@@, \uline{the canonical homomorphism} @@M8@@ \uline{is an isomorphism}. \uline{Then the homomorphism}:
@@M9@@
\uline{is surjective}. \uline{If in addition} @@M10@@ \uline{has a section} @@M11@@ \uline{of finite Tor-dimension}, \uline{then} @@M12@@ \uline{is an isomorphism}.

Take for the family @@M13@@ the family consisting of the single element @@M14@@; we then have @@M15@@, and, the conditions of 1.2 being satisfied, @@M16@@ is surjective.

If @@M17@@ has a section @@M18@@ of finite Tor-dimension, one can define a homomorphism @@M19@@ (IV 2.12). Since we have @@M20@@ by IV 2.12.1.1, @@M21@@ is injective, hence an isomorphism.

\sgaStableTarget{sga6.e09.rem.1-5}\textbf{Remark 1.5.} For @@M22@@ to have finite Tor-dimension, it is necessary and sufficient that the fibers of @@M23@@ be regular at the points of @@M24@@, i.e.\ that @@M25@@ be regular at the points of @@M26@@ (or smooth at these points if we suppose @@M27@@
